% 第 1 章 Defining Modes, Corners, and Scenarios
\bichapter{定义模式、corner 与 scenario}{Defining Modes, Corners, and Scenarios}
\label{chap:mcs}

一个 block 可能在多种条件下工作，例如不同的温度与电压，也可能处于多种功能模式。做时序分析时，每一组条件用一个 corner 表示，每一种功能模式用一个 mode 表示。scenario 是 corner 与 mode 的组合，用于时序分析与优化。

开始处理 block 之前，必须定义该 block 使用的 mode、corner 与 scenario，并定义与之关联的约束。本章包括以下主题：
\begin{itemize}
  \item 创建与删除 mode
  \item 查询 mode
  \item 创建与删除 corner
  \item 查询 corner
  \item 创建 scenario
  \item 设置 scenario 的有效分析类型
  \item 查询 scenario
  \item 删除重复的 scenario、mode 与 corner
  \item IC Compiler II 时序约束
  \item 设置时序约束
  \item 为 IC Compiler II 时序分析转换 SDC 文件
  \item 从 Design Compiler 或 IC Compiler 导入时序约束
  \item 删除 SDC 约束
\end{itemize}

\section{创建与删除 mode}
\subsection*{Creating and Removing Modes}

用 \cmd{create\_mode} 创建功能模式。每个 mode 的名称必须唯一。

例如，创建名为 \cmd{mode1} 的 mode：
\begin{lstlisting}
icc2_shell> create_mode mode1
\end{lstlisting}

用 \cmd{current\_mode} 把当前 mode 切换到指定 mode：
\begin{lstlisting}
icc2_shell> current_mode mode1
\end{lstlisting}

用 \cmd{remove\_modes} 删除当前 block 的 mode。\opt{-all} 删除全部 mode；也可使用 Tcl 列表或 collection 指定要删除的 mode。删除某个 mode 时，与该 mode 关联的全部 scenario 一并删除。

例如，删除 \cmd{mode1} 与 \cmd{mode2}：
\begin{lstlisting}
icc2_shell> remove_modes {mode1 mode2}
\end{lstlisting}

\section{查询 mode}
\subsection*{Querying Modes}

可获取以下信息：
\begin{itemize}
  \item 当前 mode：不带参数调用 \cmd{current\_mode}。
\begin{lstlisting}
icc2_shell> current_mode
{"mode1"}
\end{lstlisting}
  \item block 上已定义的 mode 列表：使用 \cmd{get\_modes}。\\
        默认返回包含当前 block 全部 mode 的 collection（与 \cmd{all\_modes} 类似）。要取得其他 block 的 mode，使用 \opt{-design}。\\
        可按属性值（\opt{-filter}）、关联的 corner（\opt{-of\_objects}）或名称模式过滤返回的 mode。
  \item 当前 block 一个或多个 mode 的详细信息：使用 \cmd{report\_modes}。\\
        报告包含每个 mode 关联的 scenario 及其有效分析类型。
\end{itemize}

报告当前 block 全部 mode 的详细信息：
\begin{lstlisting}
icc2_shell> report_modes
\end{lstlisting}

报告名为 \cmd{mode1} 与 \cmd{mode2} 的 mode：
\begin{lstlisting}
icc2_shell> report_modes {mode1 mode2}
\end{lstlisting}

\begin{seeAlsoBox}
设置 scenario 的有效分析类型
\end{seeAlsoBox}

\section{创建与删除 corner}
\subsection*{Creating and Removing Corners}

用 \cmd{create\_corner} 创建时序 corner。每个 corner 的名称必须唯一。

例如，创建名为 \cmd{corner1} 的 corner：
\begin{lstlisting}
icc2_shell> create_corner corner1
\end{lstlisting}

用 \cmd{current\_corner} 切换当前 corner：
\begin{lstlisting}
icc2_shell> current_corner corner1
\end{lstlisting}

用 \cmd{remove\_corners} 删除当前 block 的 corner。\opt{-all} 删除全部 corner；也可使用 Tcl 列表或 collection 指定 corner。删除某个 corner 时，与该 corner 关联的全部 scenario 一并删除。

例如，删除 \cmd{corner1} 与 \cmd{corner2}：
\begin{lstlisting}
icc2_shell> remove_corners {corner1 corner2}
\end{lstlisting}

\section{查询 corner}
\subsection*{Querying Corners}

可获取以下信息：
\begin{itemize}
  \item 当前 corner：不带参数调用 \cmd{current\_corner}。
\begin{lstlisting}
icc2_shell> current_corner
{"corner1"}
\end{lstlisting}
  \item block 上已定义的 corner 列表：使用 \cmd{get\_corners}。\\
        默认返回包含当前 block 全部 corner 的 collection（与 \cmd{all\_corners} 类似）。其他 block 使用 \opt{-design}。\\
        可按属性值（\opt{-filter}）、关联对象（\opt{-of\_objects}）或名称模式过滤。
  \item 当前 block 一个或多个 corner 的详细信息：使用 \cmd{report\_corners}。\\
        报告包含每个 corner 关联的 scenario 及其有效分析类型。
\end{itemize}

报告当前 block 全部 corner：
\begin{lstlisting}
icc2_shell> report_corners
\end{lstlisting}

报告名为 \cmd{corner1} 与 \cmd{corner2} 的 corner：
\begin{lstlisting}
icc2_shell> report_corners {corner1 corner2}
\end{lstlisting}

\section{创建 scenario}
\subsection*{Creating a Scenario}

用 \cmd{create\_scenario} 创建时序 scenario。每个 scenario 的名称必须唯一。默认情况下，新建 scenario 组合当前 mode 与当前 corner。用 \opt{-mode} 指定 mode，用 \opt{-corner} 指定 corner。

例如，创建名为 \cmd{scenario1} 的 scenario，组合 \cmd{mode1} 与 \cmd{corner1}：
\begin{lstlisting}
icc2_shell> create_scenario -mode mode1 -corner corner1 -name scenario1
\end{lstlisting}

创建 scenario 时，默认它处于 active 状态、对全部分析类型启用，并成为当前 scenario。同时，它的 mode 成为当前 mode，它的 corner 成为当前 corner。

用 \cmd{set\_scenario\_status -active true|false} 改变 scenario 的 active 状态。用 \cmd{set\_scenario\_status} 修改其有效分析类型。

用 \cmd{remove\_scenarios} 删除当前 block 的 scenario。\opt{-all} 删除全部 scenario；也可使用 Tcl 列表或 collection 指定要删除的 scenario。

例如，删除 \cmd{scenario1} 与 \cmd{scenario2}：
\begin{lstlisting}
icc2_shell> remove_scenarios {scenario1 scenario2}
\end{lstlisting}

\begin{seeAlsoBox}
设置 scenario 的有效分析类型
\end{seeAlsoBox}

\section{设置 scenario 的有效分析类型}
\subsection*{Setting the Active Analysis Types for a Scenario}

默认情况下，scenario 处于 active 状态，并对全部分析类型启用。用 \cmd{set\_scenario\_status} 修改一个或多个 scenario 的有效分析类型，选项见表~\ref{tab:scenario-status}。

可用 \opt{-active false} 临时关闭各项分析设置。使用该选项时，指定 scenario 的全部分析被禁用，但各项设置仍保留在设计库中。再用 \opt{-active true} 时，各项设置重新生效。

\begin{table}[htbp]
\centering
\caption{\cmd{set\_scenario\_status} 的分析选项}
\label{tab:scenario-status}
\begin{tabular}{@{}ll@{}}
\toprule
分析类型 & 命令选项 \\
\midrule
全部 & \cmd{-all} \\
无 & \cmd{-none} \\
建立时间（setup） & \cmd{-setup true | false} \\
保持时间（hold） & \cmd{-hold true | false} \\
泄漏功耗 & \cmd{-leakage\_power true | false} \\
动态功耗 & \cmd{-dynamic\_power true | false} \\
最大转换时间 & \cmd{-max\_transition true | false} \\
最大电容 & \cmd{-max\_capacitance true | false} \\
最小电容 & \cmd{-min\_capacitance true | false} \\
\bottomrule
\end{tabular}
\end{table}

\section{查询 scenario}
\subsection*{Querying Scenarios}

可获取以下信息：
\begin{itemize}
  \item 当前 scenario：不带参数调用 \cmd{current\_scenario}。
\begin{lstlisting}
icc2_shell> current_scenario
{"scenario1"}
\end{lstlisting}
  \item block 上已定义的 scenario：使用 \cmd{get\_scenarios}。\\
        默认返回包含当前 block 全部 scenario 的 collection（与 \cmd{all\_scenarios} 类似）。\\
        可按属性值（\opt{-filter}）、关联 corner（\opt{-corners}）、关联 mode（\opt{-modes}）或名称模式过滤。
  \item 当前 block 一个或多个 scenario 的详细信息：使用 \cmd{report\_scenarios}。\\
        报告包含每个 scenario 的名称、关联 mode、关联 corner 以及有效分析类型。\\
        默认报告当前 block 的全部 scenario。还可以：
  \begin{itemize}
    \item 用 \opt{-scenarios} 报告指定 scenario；
    \item 用 \opt{-modes} 报告与指定 mode 关联的 scenario；
    \item 用 \opt{-corners} 报告与指定 corner 关联的 scenario。
  \end{itemize}
\end{itemize}

报告当前 block 全部 scenario：
\begin{lstlisting}
icc2_shell> report_scenarios
\end{lstlisting}

报告名为 \cmd{scenario1} 与 \cmd{scenario2} 的 scenario：
\begin{lstlisting}
icc2_shell> report_scenarios -scenarios {scenario1 scenario2}
\end{lstlisting}

\section{删除重复的 scenario、mode 与 corner}
\subsection*{Removing Duplicate Scenarios, Modes, and Corners}

为改善运行时间与容量，可用 \cmd{remove\_duplicate\_timing\_contexts} 删除重复的 scenario、corner 与 mode。该命令会：
\begin{itemize}
  \item 删除重复且 scenario 状态相同的 scenario。
  \item 把重复 mode 上的 scenario 重新归到一个公共 mode，并删除重复 mode。重新分配 scenario 时，必要时会创建新的 corner。
  \item 把重复 corner 上的 scenario 重新归到一个公共 corner，并删除重复 corner。
\end{itemize}

判断重复 mode 时，工具不考虑已反标的开关活动（switching activity）。因此运行该命令后需要重新施加开关活动。

若使用 \cmd{set\_block\_to\_top\_map} 把 block 级的 mode、corner 与 scenario 映射到顶层，运行 \cmd{remove\_duplicate\_timing\_contexts} 之后必须手工更新该映射。用 \cmd{report\_block\_to\_top\_map} 查看顶层与 block 级 mode、corner、scenario 之间的关系。

\section{IC Compiler II 时序约束}
\subsection*{IC Compiler II Timing Constraints}

IC Compiler II 把时序约束分为以下几类：
\begin{itemize}
  \item \textbf{Mode 专用约束}\\
        定义或修改时序图（timing graph，表示电路时序的加权图）的约束是 mode 专用约束，包括：
  \begin{itemize}
    \item 时钟定义
    \item 最大与最小延迟约束
  \end{itemize}
        至少应为每个时钟信号包含一条时钟定义。
  \item \textbf{Corner 专用约束}\\
        修改对象上计算延迟的约束是 corner 专用约束，包括：
  \begin{itemize}
    \item 工作条件（operating conditions）
    \item 端口负载电容
    \item 寄生参数信息与缩放因子
    \item 时序 derating 因子
  \end{itemize}
  \item \textbf{Scenario 专用约束}\\
        带有延迟值并引用 modal 对象（例如时钟）的约束是 scenario 专用约束，包括：
  \begin{itemize}
    \item 输入与输出延迟
    \item 时钟特性，例如 clock latency、uncertainty 与 transition
  \end{itemize}
        至少应为每个 I/O 端口包含输入延迟或输出延迟。
  \item \textbf{网表专用约束}\\
        适用于全部 corner、mode 与 scenario 的约束是网表专用约束，包括 ideal network 设置。\\
        它们可以全局作用于全部设计，也可以局部作用于某个设计。例如 \cmd{set\_ideal\_network} 设置的约束就是网表专用的。
\end{itemize}

\section{设置时序约束}
\subsection*{Setting Timing Constraints}

可以通过读入 Synopsys Design Constraints（SDC）文件，或逐条执行 SDC 命令来指定时序约束。SDC 命令的细节见 \textit{Using the Synopsys Design Constraints Format Application Note}。

用 \cmd{read\_sdc} 读入 SDC 文件：
\begin{lstlisting}
icc2_shell> read_sdc block.sdc
\end{lstlisting}

\begin{noteBox}
若 SDC 文件不含单位设置，单位取自主逻辑库。若 SDC 文件含有单位设置，则必须与主逻辑库中的单位一致。
\end{noteBox}

为使 block 设置高效且准确，按以下步骤操作：
\begin{enumerate}
  \item 把 block 约束拆成下列独立文件：
  \begin{itemize}
    \item 每个 mode 一个 mode 专用文件
    \item 每个 corner 一个 corner 专用文件
    \item 每个 scenario 一个 scenario 专用文件
    \item block 的全局约束文件
  \end{itemize}
  \item 为该 block 定义全部 mode、corner 与 scenario。
  \item 在施加对应约束之前，先把当前 mode、corner 或 scenario 设对。\\
        读入 SDC 文件或执行 SDC 命令时，约束按类型只作用于当前 scenario，或其关联的 mode / corner。
\end{enumerate}

例如，假设某 block 具有：
\begin{itemize}
  \item 两个 mode：\cmd{M1} 与 \cmd{M2}
  \item 一个 corner：\cmd{C}
  \item 两个 scenario：
  \begin{itemize}
    \item \cmd{M1@C}，由 mode \cmd{M1} 与 corner \cmd{C} 组合
    \item \cmd{M2@C}，由 mode \cmd{M2} 与 corner \cmd{C} 组合
  \end{itemize}
\end{itemize}

下列脚本创建 mode、corner 与 scenario，并施加对应约束与设置：
\begin{lstlisting}
# Create all modes, corners, and scenarios
create_mode M1
create_mode M2
create_corner C
create_scenario -mode M1 -corner C -name M1@C
create_scenario -mode M2 -corner C -name M2@C
# Apply mode-, corner-, and scenario-specific constraints
current_mode M1
read_sdc M1_mode.sdc
current_mode M2
read_sdc M2_mode.sdc
current_corner C
read_sdc C_corner.sdc
## Apply TLUPlus settings, extraction options, and other settings
current_scenario M1@C
read_sdc M1@C_scenario.sdc
current_scenario M2@C
read_sdc M2@C_scenario.sdc
\end{lstlisting}

\section{为 IC Compiler II 时序分析转换 SDC 文件}
\subsection*{Converting an SDC File for IC Compiler II Timing Analysis}

IC Compiler II 只支持 OCV 时序分析。若 SDC 按 best-case/worst-case（BCWC）时序分析编写，必须把它转换成两份 OCV SDC 文件。

把 BCWC SDC 转换成两份 OCV SDC 的步骤：
\begin{enumerate}
  \item 用 \cmd{create\_mode} 为设计创建一个 mode，见「创建与删除 mode」。
\begin{lstlisting}
icc2_shell> create_mode my_mode
\end{lstlisting}
  \item 为 best-case 约束创建：
  \begin{itemize}
    \item 用 \cmd{create\_corner} 创建一个 corner；
    \item 用 \cmd{create\_scenario} 创建一个 scenario。
  \end{itemize}
        best-case scenario 不用于 setup 分析，因此用 \cmd{set\_scenario\_status -setup false} 关闭该 scenario 的 setup 分析。
\begin{lstlisting}
icc2_shell> create_corner bc_corner
icc2_shell> create_scenario -mode my_mode -corner bc_corner \
    -name bc_scenario
icc2_shell> set_scenario_status -setup false bc_scenario
\end{lstlisting}
  \item 用 \cmd{read\_sdc} 从 BCWC SDC 文件读入 best-case 约束。\\
        默认情况下，工具读入 SDC 文件中的全部约束。若只保留 best-case 约束，读入前把应用选项 \appopt{time.convert\_constraint\_from\_bc\_wc} 设为 \cmd{bc\_only}。
\begin{lstlisting}
icc2_shell> set_app_options \
    -name time.convert_constraint_from_bc_wc -value bc_only
icc2_shell> read_sdc bc_wc.sdc
\end{lstlisting}
  \item 为 worst-case 约束创建 corner 与 scenario。\\
        worst-case scenario 不用于 hold 分析，因此用 \cmd{set\_scenario\_status -hold false} 关闭 hold 分析。
\begin{lstlisting}
icc2_shell> create_corner wc_corner
icc2_shell> create_scenario -mode my_mode -corner wc_corner \
    -name wc_scenario
icc2_shell> set_scenario_status -hold false wc_scenario
\end{lstlisting}
  \item 用 \cmd{read\_sdc} 从同一 BCWC SDC 文件读入 worst-case 约束。\\
        读入前把 \appopt{time.convert\_constraint\_from\_bc\_wc} 设为 \cmd{wc\_only}。
\begin{lstlisting}
icc2_shell> set_app_options \
    -name time.convert_constraint_from_bc_wc -value wc_only
icc2_shell> read_sdc bc_wc.sdc
\end{lstlisting}
  \item 把 \appopt{time.convert\_constraint\_from\_bc\_wc} 复位为 \cmd{none}。
\begin{lstlisting}
icc2_shell> set_app_options \
    -name time.convert_constraint_from_bc_wc -value none
\end{lstlisting}
\end{enumerate}

\section{从 Design Compiler 或 IC Compiler 导入时序约束}
\subsection*{Importing Timing Constraints From the Design Compiler or IC Compiler Tool}

把 Design Compiler 或 IC Compiler 的时序约束导入 IC Compiler II：
\begin{enumerate}
  \item 在 Design Compiler 或 IC Compiler 中用 \cmd{write\_timing\_context -format icc2} 生成约束。\\
        必须用 \opt{-output} 指定输出名称。输出是一个目录，其中包含表~\ref{tab:timing-context} 所列 Tcl 脚本。
\end{enumerate}

\begin{table}[htbp]
\centering
\caption{\cmd{write\_timing\_context -format icc2} 生成的 Tcl 脚本}
\label{tab:timing-context}
\begin{tabularx}{\textwidth}{@{}l X@{}}
\toprule
脚本名 & 内容 \\
\midrule
\cmd{top.tcl} &
创建 mode、corner 与 scenario 的命令；以及向 mode、corner、scenario 和设计填入约束的命令 \\
\cmd{scenario\_<scenario\_name>.tcl} & 特定 scenario 的约束 \\
\cmd{design.tcl} & 与 scenario 无关的设计约束 \\
\bottomrule
\end{tabularx}
\end{table}

默认情况下，该命令为全部 scenario（active 与 inactive）以及 block 上与 scenario 无关的时序约束生成文件。若只生成指定 scenario 的约束，使用 \opt{-scenarios}。

\begin{enumerate}
  \setcounter{enumi}{1}
  \item 在 IC Compiler II 中 source \cmd{top.tcl}，从而创建对应的 mode、corner 与 scenario 并施加约束。
\end{enumerate}

下列示例在 Design Compiler 中为全部 active scenario 以及与 scenario 无关的时序约束生成文件：
\begin{lstlisting}
dc_shell> write_timing_context -scenarios [get_scenarios -active true] \
    -output BLK1_const
\end{lstlisting}

下列示例把 Design Compiler 生成的约束导入 IC Compiler II：
\begin{lstlisting}
icc2_shell> open_block BLK1
icc2_shell> source BLK1_const/top.tcl
\end{lstlisting}

\section{删除 SDC 约束}
\subsection*{Removing SDC Constraints}

用 \cmd{remove\_sdc} 从 block 中删除 SDC 约束。

默认删除当前 block 中的全部 SDC 约束。也可以这样控制删除范围：
\begin{itemize}
  \item 用 \opt{-modes}、\opt{-corners} 或 \opt{-scenarios} 删除指定 mode、corner 或 scenario 上的约束。\\
        例如，删除 mode \cmd{M1} 的 mode 专用约束、corner \cmd{C1} 与 \cmd{C2} 的 corner 专用约束，以及 scenario \cmd{M1C1} 与 \cmd{M1C2} 的 scenario 专用约束：
\begin{lstlisting}
icc2_shell> remove_sdc -modes M1 -corners {C1 C2} \
    -scenarios {M1C1 M1C2}
\end{lstlisting}
  \item 用 \opt{-design} 删除不关联任何 mode、corner 或 scenario 的设计专用约束。
  \item 用 \opt{-include} 或 \opt{-exclude} 包含或排除特定类型的约束。
\end{itemize}

只删除 case analysis 设置：
\begin{lstlisting}
icc2_shell> remove_sdc -include case_analysis
\end{lstlisting}

删除除 ideal network 设置以外的全部 SDC 约束：
\begin{lstlisting}
icc2_shell> remove_sdc -exclude ideal_network
\end{lstlisting}

\opt{-include} 与 \opt{-exclude} 的完整合法取值见 \cmd{remove\_sdc} 的 man 页。
